% Figura: diagrama de flujo del método
\begin{figure}[!htb]
\centering
\begin{tikzpicture}[node distance=0.55cm]
  \tikzset{paso/.style={nodo, text width=7.4cm, align=left, minimum height=0.85cm},
           num/.style={font=\sffamily\bfseries\small, text=ambar}}
  \node[nodocian, align=center, text width=5.4cm] (ini)
    {$P(x)\,y''+A(x)\,y'+B(x)\,y=0$\\[1pt]\footnotesize\color{gris} y un punto $x_0$};
  \node[diamond, aspect=2.4, draw=violeta, fill=cajafondo, line width=0.6pt,
        inner sep=1.5pt, font=\small, below=of ini] (dec) {¿$P(x_0)\neq0$?};
  \node[paso, below=0.75cm of dec] (p1)
    {\textcolor{ambar}{\bfseries 1.}\ \ $R\ge\dist\bigl(x_0,\mathcal Z(P)\bigr)$\quad(ceros complejos)};
  \node[paso, below=of p1] (p2)
    {\textcolor{ambar}{\bfseries 2.}\ \ $y=\sum a_n z^{n}$; derivar término a término};
  \node[paso, below=of p2] (p3)
    {\textcolor{ambar}{\bfseries 3.}\ \ desplazar índices y producto de Cauchy};
  \node[paso, below=of p3] (p4)
    {\textcolor{ambar}{\bfseries 4.}\ \ coeficiente de $z^{n}$ nulo $\Rightarrow$ recurrencia};
  \node[paso, below=of p4] (p5)
    {\textcolor{ambar}{\bfseries 5.}\ \ $(1,0)\mapsto y_1$, \ $(0,1)\mapsto y_2$; \ $y=a_0y_1+a_1y_2$};
  \node[nodo, draw=rojo, align=center, text width=3.6cm, right=1.6cm of dec] (sing)
    {\textcolor{rojo}{\bfseries Punto singular}\\[1pt]
     \footnotesize\color{gris} serie de Frobenius\\(Parte II)};

  \draw[flecha] (ini) -- (dec);
  \draw[flecha] (dec) -- node[right=2pt, etiqgris] {sí} (p1);
  \draw[flecha] (dec) -- node[above=2pt, etiqgris] {no} (sing);
  \foreach \a/\b in {p1/p2, p2/p3, p3/p4, p4/p5}{\draw[flecha] (\a) -- (\b);}
\end{tikzpicture}
\caption{Diagrama de flujo del método de la serie entera.}
\label{fig:flujo}
\end{figure}
